% META: {"id":"1NSI-TC-CO-027","chapitre":"1NSI-TYPES-CONSTRUITS","type_objet":"corrige","capacites":["1NSI-TYPES-CONSTRUITS-C3"],"mode_creation":"ex_nihilo","sources_inspiration":[],"status":"approved","exercice_ref":"1NSI-TC-EX-027","origin":"needs_review","status_history":["needs_review","audited","accepted","approved"],"release_acceptance":"RELEASE_OWNER_BATCH_ACCEPTANCE","acceptance_closure_digest":"sha256:9b3ccf9a81c5520fbb7e03b7b2d3e7bf2057a02834b6d908bf3be5f49f3b553f","acceptance_receipt_digest":"sha256:4fad86f0282e0fa4e0419cca1ad6379ae7af5a280c04a4a90880f90a1c044242"}
% BEGIN-VERIFY
% def creer_grille(n, p, valeur):
%     grille = []
%     for i in range(n):
%         ligne = []
%         for j in range(p):
%             ligne.append(valeur)
%         grille.append(ligne)
%     return grille
% g = creer_grille(2, 3, 0)
% assert g == [[0, 0, 0], [0, 0, 0]]
% assert len(g) == 2
% assert len(g[0]) == 3
% g[0][0] = 5
% assert g[1][0] == 0
% END-VERIFY
\begin{corrige}{1NSI-TC-EX-027}
\begin{python}
def creer_grille(n, p, valeur):
    grille = []
    for i in range(n):
        ligne = []
        for j in range(p):
            ligne.append(valeur)
        grille.append(ligne)
    return grille
\end{python}

L'écriture \lstinline{[[valeur] * p] * n} est incorrecte car l'opérateur \lstinline{*} sur une liste crée \lstinline{n} références vers \emph{la même} sous-liste. Modifier une case d'une ligne modifierait toutes les lignes (effet de bord dû à l'aliasing).
\end{corrige}
